\documentclass[12pt]{article}

\input{../../_assets/preambulo.tex}


\begin{document}

    % 1. Foto de fondo
    % 2. Título
    % 3. Encabezado Izquierdo
    % 4. Color de fondo
    % 5. Coord x del titulo
    % 6. Coord y del titulo
    % 7. Fecha

    
    \input{../../_assets/portada}
    \portadaExamen{ffccA4.jpg}{Álgebra II\\Examen X}{Álgebra II. Examen X}{MidnightBlue}{-8}{28}{2026}{David Muñoz Gómez}

    \begin{description}
        \item[Curso Académico] 25/26.
        \item[Grado] DGIIM.
        \item[Grupo] Único.
        \item[Profesor] Aurora Inés del Río Cabeza.
        \item[Descripción] Convocatoria Ordinaria.
        \item[Fecha] 17/06/2026.
        % \item[Duración] ---.
    \end{description}
    \newpage


    % ------------------------------------
    
    \begin{ejercicio}[$1.75$ puntos]
        Se considera el grupo diédrico $$D_5 = \langle r, s \mid r^5 = 1, s^2 = 1, sr = r^{-1}s \rangle$$ y el grafo $G$ cuyos vértices son los elementos de $D_5$ y en el que, para cualquier $a \in D_5$, hay una arista entre $a$ y $ra$ y también una arista entre $a$ y $sa$.
        
        \begin{enumerate}[label=\alph*)]
            \item ($0.25$ puntos) Calcula la sucesión de grados de $G$ y calcula su matriz de adyacencia.
            \item ($0.75$ puntos) Razona si $G$ es un grafo de Euler, de Hamilton o plano.
            \item ($0.75$ puntos) Considera un nuevo grafo $G'$ obtenido añadiendo a $G$ un nuevo vértice adyacente a todos los de $G$. Razona si $G'$ es un grafo de Euler o si tiene un camino de Euler. ¿Es Hamiltoniano? ¿Y plano?
        \end{enumerate}
    \end{ejercicio}

    \begin{ejercicio}[$0.75$ puntos]
        Demuestra que un grupo finito $G$ tiene un único $p$-subgrupo de Sylow para cada primo $p$ si, y solo si, $G$ es producto directo de sus subgrupos de Sylow.
    \end{ejercicio}

    \begin{ejercicio}[$1$ punto]~
        \begin{enumerate}[label=\alph*)]
            \item ($0.5$ puntos) Demuestra que no existen grupos simples de orden 351.
            \item ($0.5$ puntos) Si $G$ es un grupo finito con un subgrupo normal $N$ de índice 351, demuestra que, si $N$ es resoluble, entonces $G$ también es resoluble.
        \end{enumerate}
    \end{ejercicio}

    \begin{ejercicio}[$2.5$ puntos]
        Sea $G$ un grupo de orden 20.
        
        \begin{enumerate}[label=\alph*)]
            \item ($0.5$ puntos) Razona que $G$ es el producto semidirecto un 2 subgrupo de Sylow de $G$ y un 5 subgrupo de Sylow de $G$.
            \item ($1$ punto) Razona que si $G$ tiene un elemento de orden 4 entonces hay al menos tres productos semidirectos (no isomorfos) $P_5 \rtimes P_2$ uno de ellos abeliano y dos no abelianos, da una presentación de estos últimos.
            \item ($1$ punto) Concluye que hay al menos 5 grupos de orden 20, dos abelianos y tres no abelianos. Da las descomposiciones cíclicas de los abelianos y presentaciones de los no abelianos.
        \end{enumerate}
    \end{ejercicio}

    \begin{ejercicio}[$2$ puntos]
        Sea $A = \bb{Z}^4/K$ donde $K$ es el subgrupo con generadores
        \[ \{(5, 90, 130, 0), (-25, 120, 25, -30), (5, 0, 10, 0)\}. \]
        
        \begin{enumerate}[label=\alph*)]
            \item ($1.25$ puntos) Calcula, de forma razonada, el rango (de la parte libre) y las descomposiciones cíclica y cíclica primaria de $A$. Describe cuales son las componentes $p$-primarias no nulas de $A$ y razona si $A$ tiene elementos de orden $\infty$, de orden 50 o de orden 6 dando, en su caso, un ejemplo de cada uno de ellos.
            \item ($0.75$ puntos) Clasifica, dando sus descomposiciones cíclica y cíclica primaria, todos los grupos abelianos cuyo orden sea igual al orden de $T(A)$, el grupo de torsión de $A$. Determina la longitud y los factores de composición de todos ellos.
        \end{enumerate}
    \end{ejercicio}

    \begin{ejercicio}[$2$ puntos]
        Define el concepto de grupo simple. Demuestra el Teorema de Abel.
    \end{ejercicio}

    \newpage
    \setcounter{ejercicio}{0} % Reiniciar contador de ejercicios
    \noindent
    % \textbf{Solución.}

\end{document}